\chapter{Paradoxos de Medida e Geometria}
\label{chap:apendiceMedidaGeometria}

Os paradoxos deste apêndice habitam o espaço geométrico e revelam que o Axioma da Escolha, discutido no \autoref{chap:escolha} tem consequências que violam profundamente a intuição física sobre volume, medida e decomposição.

\section{O Conjunto de Vitali (1905)}

Apresentar a construção de Vitali: usando o Axioma da Escolha,
escolhe-se um representante de cada classe de equivalência de \(\mathbb{R}-\mathbb{Q}\)
no intervalo \([0,1]\). O conjunto resultante não possui medida de Lebesgue
bem definida.
Situar como o antecessor direto e necessário de Banach-Tarski:
ambos dependem essencialmente do Axioma da Escolha e de conjuntos
não mensuráveis.

\section{O Conjunto de Cantor (1883)}

Apresentar a construção por remoções sucessivas de terços centrais
do intervalo \([0,1]\). Demonstrar que o conjunto resultante tem
medida de Lebesgue zero (soma das remoções igual a~1) e ainda
assim cardinalidade \(\mathfrak{c}\), via representação ternária.
Contrastar diretamente com C.1 (Vitali): enquanto o conjunto de
Vitali exige o Axioma da Escolha para ser construído, o Conjunto
de Cantor é inteiramente explícito, e mesmo assim produz a mesma
tensão entre medida e cardinalidade. Conectar com o Capítulo~6:
assim como dimensão e cardinalidade são independentes, medida e
cardinalidade também o são.

\section{A Lâmpada de Thomson (1954)}

Apresentar o experimento: uma lâmpada é ligada e desligada infinitas
vezes em tempo finito (em t = 1/2, 1/4, 1/8, ...).
Qual é seu estado em t = 0?
Discutir como operações infinitas nem sempre convergem a um estado
bem definido.


\section{Os Jogos de Banach-Mazur (1935)}

Apresentar o jogo de dois jogadores sobre subconjuntos de reais:
um jogador tenta ``cair'' em um conjunto alvo, o outro tenta evitar.
Mostrar que a existência de estratégias vencedoras depende da
complexidade descritiva do conjunto alvo.
Mencionar que o Axioma da Determinação (incompatível com o Axioma
da Escolha) tornaria todo jogo desse tipo determinado.
Conectar com o Capítulo 8 como extensão natural da discussão sobre
o Axioma da Escolha.
